\section{Alimentación seleccionada y vigilancia de batería (INC-42)}
\label{sec:sd4-inc42}
Se conservan las cuatro FRPS y los cuatro VEP ordinarios de INC-39, pero se selecciona montaje adyacente fuente/detector, dentro de la misma zona F01/F02/F03/F06. Cada ramal DC completo, incluidas holguras de conexión, queda limitado a un metro de ida. Este límite sustituye los cinco metros anteriores. Se selecciona cable Belden \texttt{5100UL 0021000}, rojo, FPLR, dos conductores AWG14 de cobre multifilar 19x27, aislamiento PP y cubierta PVC. La ficha publica diámetro nominal 0,218 pulgadas (5,5372 mm), tensión 300 V y radio mínimo 44,2 mm. El dato resistivo 2,72 $\Omega$/1000 pies es nominal, no máximo garantizado. El uso seleccionado es interior ordinario; no se atribuyen prestaciones LSZH, marinas o de zona de gas \parencite[pp.~1--2]{lote42-belden5100ul}.

El manual FRPS permite bornes 14--22 AWG, cable FPLR y conexión multifilar; VEP admite 24--14 AWG (0,2--2,5 mm$^2$). Se mantienen segregación de potencia limitada/no limitada, ferrita exigida a la salida de envolvente, alivio de tracción y terminaciones OEM. No se apilan conductores ni se puentea una borna para simular continuidad. El EOLR-1 queda conectado en el extremo VEP, con sus 20 mA y monitor independiente ya presupuestados. Se conservan 56 monitores, 128 direcciones y reparto SLC 40/40/24/24, sin nuevos módulos por este cambio \parencite[secs.~1.3, 2.1.2 y 2.5]{lote39-frpsmanual}\parencite[p.~2]{lote30-vep36106}.

La FRPS identifica batería entre 18,4 y 19,4 V mediante estado D37 y borra ese estado al alcanzar 22 V; por debajo de 18,4 V retira alimentación. Se mantiene P3 Trouble como señal independiente de fuente incluso durante alarma. Estos valores corresponden a vigilancia de batería, no a un sensor calibrado de 18 V en bornes VEP. La recepción debe correlacionar D37, P3, EOLR y Fault VEP durante descarga de batería, registrar tolerancias y tensión real en detector, y verificar que una avería de potencia no se interprete como fuego. EOLR puede permanecer activado a menos de 18 V, por lo que no demuestra detector operativo \parencite[secs.~1.3--1.4]{lote39-frpsmanual}\parencite{lote39-eolr1}.

\section{Margen de tensión y descarte de reemplazo incompatible (EXT-42)}
\label{sec:sd4-ext42}
Para el ramal seleccionado se conserva, como envolvente propia de cribado, 10,7 $\Omega$/km por conductor a 20 grados C, factor 1,13755 a 55 grados C, corriente externa total limitada a 1 A y reserva 0,10 V de conexiones. La resistencia conservadora supera los 8,923885 $\Omega$/km nominales publicados para Belden, pero no certifica el máximo de cada suministro. Resulta $\Delta V\leq2(0,001)(10,7)(1,13755)(1)+0,10=0,124344$ V, reducción de 0,097374 V frente al ramal de cinco metros. La fuente debe entregar al menos $18+0,124344=18,124344$ V en su salida para conservar 18 V en detector. Con 18,6 V reales al origen quedarían 18,475656 V. Si se comparara batería de 18,4 V con salida, la caída interna disponible sería sólo 0,275656 V; el manual no acredita ese máximo interno ni su tolerancia de corte. No se equipara tensión de batería con salida AUX.

Se mantienen 0,625/0,751 A por FRPS a 24 V, cribado 18,0751 Ah, pares de 33 Ah y presupuestos de Cheetah/SPS de EXT-39. El máximo VEP publicado a 24 V no demuestra corriente bajo 1 A a 18 V; se requiere curva OEM o caracterización aceptada del conjunto y capacidad útil del par a temperatura y tensión final. Acortar cable mejora margen, pero no completa autonomía ni recarga. La liberación DFIA conserva su estudio independiente: este ramal no justifica sus 22 V finales, ni concentración, hidráulica, retención o gas residual.

Se examinó un reemplazo OEM concreto: VPS-100US-UL/VPS-300US-UL Gen 2. Su hoja de instalación identifica VEP compatible, listado UL 864, salida de aplicación especial 19,96--27,54 V y hasta 1,2 A por circuito/1,5 A total. Con la envolvente de este ramal, 19,96 V produciría 19,835656 V al extremo. Sin embargo, el OEM declara entrada 120/220 V a 60 Hz; la alimentación de esta instalación es 50 Hz. Se descarta su incorporación actual sin aprobación explícita de una variante compatible. Tampoco se adopta VPS-250-E5 como reemplazo UL: su ficha identifica EN54-4/VdS, salida 19,5--30 V y 230 V AC, sin declarar ese listado UL. No se incorpora un convertidor genérico entre fuente y VEP \parencite{lote42-vpsinst,lote42-e5}.

El manual Gen 2 permite seis baterías de 12 Ah o un par de 36 Ah para ciertos perfiles de carga, exige gabinete VBC-001 para expansión y remite al calculador 21062. Un par de 12 Ah no cubre siquiera el subtotal VEP/EOLR a 24 V: $1,2[24(0,590)+(5/60)(0,610)]=17,053$ Ah, antes del consumo propio de esa fuente. Su tabla de 500 mA tampoco representa el VEP elegido. El contacto de batería baja no publica umbral/tolerancia; el relé remoto retrasa la pérdida AC, mientras el local la comunica inmediatamente. No se traslada esa lógica ni baterías VBT al inventario FRPS. Estos descartes evitan atribuir cierre a una alternativa aún incompatible \parencite{lote42-vpsinst}.

\section{Ramal interior y conservación de pasos (PEN-42)}
\label{sec:sd4-pen42}
Los cuatro ramales INC-F01/F02/F03/F06-PWR/EOL permanecen identificados. Se fija el tramo de potencia fuente/VEP íntegramente dentro de su zona ordinaria, en conducto propio y sin cruzar barrera de compartimentación. Son cuatro tramos de hasta un metro, máximo cuatro metros instalados de cable, sin sumar holguras después del límite. La pareja de monitorización de contacto conserva su ruta y selección pendientes; no se declara que el nuevo cable de potencia resuelva todos los servicios de control.

Se conservan treinta aberturas y sesenta conjuntos SF/Wedge por los demás servicios. El diámetro nominal del cable seleccionado sirve para ordenar instalación interior, sin convertir su cubierta PVC/aislamiento PP en aprobación C-AJ-8309. Si una disposición exige atravesar barrera, deja de corresponder al tramo interior seleccionado y deberá resolver compatibilidad de aislamiento, diámetro real y empaquetado antes de recibir ese paso. El cambio no elimina obligaciones de cables restantes, piso, juntas, soportes o enlaces marinos. RT-06.04, RT-06.16 y RT-06.17 conservan brechas y estado pendiente.
